\documentclass[10pt]{article}

\usepackage[margin=0.95in]{geometry}
\usepackage[T1]{fontenc}
\usepackage{lmodern}
\usepackage{amsmath,amssymb,amsthm,mathtools}
\usepackage[scaled=0.95]{inconsolata}
\usepackage{microtype}
\usepackage{xcolor}
\usepackage[colorlinks=true,linkcolor=blue!55!black,citecolor=blue!55!black,urlcolor=blue!55!black]{hyperref}
\usepackage[nameinlink,noabbrev]{cleveref}

\numberwithin{equation}{section}
\setlength{\abovedisplayskip}{6pt}
\setlength{\belowdisplayskip}{6pt}
\setlength{\abovedisplayshortskip}{4pt}
\setlength{\belowdisplayshortskip}{4pt}

\newcommand{\TuRBO}{\textnormal{\texttt{TuRBO}}}
\newcommand{\TurboOne}{\textnormal{\texttt{Turbo1}}}
\newcommand{\TurboM}{\textnormal{\texttt{TurboM}}}
\newcommand{\X}{\mathcal{X}}
\newcommand{\OmegaD}{\X_D}
\newcommand{\Dpos}{\mathcal{D}_{++}^{D}}
\newcommand{\R}{\mathbb{R}}
\newcommand{\N}{\mathbb{N}}
\newcommand{\TR}{\mathcal{T}}
\newcommand{\failtol}{\tau_{\mathrm{fail}}}
\newcommand{\best}{f_{\mathrm{best}}}
\newcommand{\indic}[1]{\mathbf{1}\!\left\{#1\right\}}
\DeclareMathOperator{\diag}{diag}
\DeclareMathOperator{\supp}{supp}
\DeclareMathOperator{\Bern}{Bernoulli}
\DeclareMathOperator{\Bin}{Binomial}
\DeclareMathOperator{\Unif}{Unif}
\DeclareMathOperator{\Vol}{Vol}
\DeclareMathOperator{\GM}{GM}
\DeclarePairedDelimiter{\norm}{\lVert}{\rVert}
\DeclarePairedDelimiter{\abs}{\lvert}{\rvert}
\DeclarePairedDelimiter{\card}{\lvert}{\rvert}
\DeclarePairedDelimiter{\floor}{\lfloor}{\rfloor}
\DeclarePairedDelimiter{\ceil}{\lceil}{\rceil}

\theoremstyle{plain}
\newtheorem{theorem}{Theorem}
\newtheorem{lemma}{Lemma}
\newtheorem{corollary}{Corollary}

\theoremstyle{definition}
\newtheorem{definition}{Definition}
\newtheorem{example}{Example}

\theoremstyle{remark}
\newtheorem{remark}{Remark}

\title{\vspace{-2.2em}Failure Tolerance in High-Dimensional \TuRBO{}}
\author{}
\date{\vspace{-1.2em}\today\vspace{-1.2em}}

\begin{document}
\maketitle

\begin{abstract}
Trust-region Bayesian optimization methods maintain a center \(x_k\), a local
region \(\TR_k\), and a length parameter \(\Delta_k\).  The length parameter is
shrunk when the local failure counter reaches a prescribed tolerance.  In the
reference implementation of \TuRBO{} \cite{eriksson2019turbo,turboCode}, the
number \(K_{k,c}\) of active RAASP mask coordinates satisfies
\[
  \mathbb{E}_D\!\left[K_{k,c}\right]=20+e^{-20}+O(D^{-1}),
  \qquad D\to\infty .
\]
The failure tolerance is ambient-dimensional: a shrink can be triggered only
after at least \(D\) evaluations are assigned through unsuccessful updates in
the same region.  We formulate this scale separation in trust-region notation
and prove
\[
  N_{\downarrow}(T)\le \floor*{\frac{T}{D}},
\]
where \(N_{\downarrow}(T)\) is the number of possible shrink events after
\(T\) assigned evaluations.  Hence \(T<D\) implies
\(N_{\downarrow}(T)=0\), even along an all-failure trajectory.
\end{abstract}

\section{Trust-Region State}

Let
\[
  \OmegaD=[0,1]^D,
  \qquad
  \Dpos=\{\,\diag(w_1,\ldots,w_D):w_j>0\,\}.
\]
At local iteration \(k\), write the trust-region state as
\begin{equation}
  \mathsf S_k=(x_k,\Delta_k,W_k)
  \in
  \OmegaD\times(0,\infty)\times\Dpos .
  \label{eq:tr-state}
\end{equation}
The scalar \(\Delta_k\) is the update variable; \(W_k\) fixes relative
coordinate scales.  The induced weighted \(\ell_\infty\) region is
\begin{equation}
  \TR_k
  =
  \TR(\mathsf S_k)
  =
  \left\{
    x\in\OmegaD:
    \norm*{W_k^{-1}(x-x_k)}_\infty
    \le
    \frac{\Delta_k}{2}
  \right\}.
  \label{eq:trust-region}
\end{equation}
This is the standard trust-region object, specialized to the axis-aligned
hyperrectangles used by \TuRBO{} \cite{conn2000trust,eriksson2019turbo}.

In \TuRBO{}, \(W_k\) is determined by local GP lengthscales.  For
\(\ell_k^{\mathrm{GP}}\in(0,\infty)^D\), define
\begin{equation}
  g_k
  =
  \left(\prod_{h=1}^{D}\ell_{k,h}^{\mathrm{GP}}\right)^{\frac{1}{D}},
  \qquad
  w_{k,j}
  =
  \frac{\ell_{k,j}^{\mathrm{GP}}}{g_k},
  \qquad
  W_k=\diag(w_{k,1},\ldots,w_{k,D}).
  \label{eq:shape-matrix}
\end{equation}
Then
\begin{equation}
  \det W_k=1,
  \qquad
  L_{k,j}=\Delta_k w_{k,j},
  \qquad
  \GM(L_{k,1},\ldots,L_{k,D})=\Delta_k .
  \label{eq:side-lengths}
\end{equation}
Away from boundary clipping, \(\Vol(\TR_k)=\prod_{j=1}^{D}L_{k,j}=\Delta_k^D\).
ENN-surrogate \TuRBO{} uses the isotropic specialization \(W_k=I_D\).  Hence
\(L_{k,j}=\Delta_k\) for every coordinate.

\section{RAASP Candidate Support}

\TuRBO{} generates a candidate cloud inside \(\TR_k\) by applying a random
axis-aligned subspace perturbation (RAASP) to a base trust-region displacement.
Write
\[
  [D]=\{1,\ldots,D\},
  \qquad
  e_j=(0,\ldots,0,1,0,\ldots,0)\in\{0,1\}^{D}.
\]
For candidate \(c\), let
\begin{align}
  u_{k,c}
  &\in
  \left[-\frac{\Delta_k}{2},\frac{\Delta_k}{2}\right]^D,
  \label{eq:base-displacement}
  \\
  \bar b_{k,c,j}
  &\stackrel{\mathrm{ind}}{\sim}\Bern(p_D),
  \qquad j\in[D],
  \label{eq:raw-mask}
  \\
  p_D
  &=
  \min\left\{\frac{20}{D},1\right\}.
  \label{eq:mask-prob}
\end{align}
Set
\begin{equation}
  \bar Z_{k,c}=\sum_{j=1}^{D}\bar b_{k,c,j}.
  \label{eq:raw-mask-size}
\end{equation}
The reference implementation then enforces a nonzero mask.  Equivalently, for
an independent \(J_{k,c}\sim\Unif([D])\),
\begin{equation}
  b_{k,c}
  =
  \begin{cases}
    \bar b_{k,c}, & \bar Z_{k,c}>0,\\
    e_{J_{k,c}}, & \bar Z_{k,c}=0.
  \end{cases}
  \label{eq:corrected-mask}
\end{equation}
This is the RAASP correction used in the reference code \cite{turboOne}.
Let \(\mathbb{P}_D\) denote the joint law of \((\bar b_{k,c},J_{k,c})\), and let
\(\mathbb{E}_D\) denote expectation under \(\mathbb{P}_D\).  The proposed and realized steps
below are conditional on the trust-region state \(\mathsf S_k\).
The proposed and realized steps are
\begin{align}
  \tilde s_{k,c}
  &=
  W_k(b_{k,c}\odot u_{k,c}),
  \label{eq:proposed-step}
  \\
  x_{k,c}
  &=
  \Pi_{\OmegaD}(x_k+\tilde s_{k,c}),
  \qquad
  s_{k,c}=x_{k,c}-x_k .
  \label{eq:realized-step}
\end{align}
Thus \(\tilde s_{k,c}\) is the proposed step before clipping, and
\(s_{k,c}\) is the realized step after projection onto \(\OmegaD\).  Boundary
clipping can only remove active coordinates:
\begin{equation}
  \supp(s_{k,c})\subseteq\supp(\tilde s_{k,c}),
  \qquad
  \norm{s_{k,c}}_0\le\norm{\tilde s_{k,c}}_0 .
  \label{eq:clipping-support}
\end{equation}

For \(s\in\R^D\), define
\begin{equation}
  A(s)=\supp(s)=\{j:s_j\ne0\},
  \qquad
  \norm{s}_0=\card{A(s)}.
  \label{eq:active-set}
\end{equation}
The mask support and its scale are
\begin{equation}
  K_{k,c}=\norm{b_{k,c}}_0,
  \qquad
  \kappa_D=\mathbb{E}_D\!\left[K_{k,c}\right].
  \label{eq:kappa-def}
\end{equation}
If \(u_{k,c,j}\ne0\) almost surely for every \(j\), then
\begin{equation}
  A(\tilde s_{k,c})=\{j:b_{k,c,j}=1\},
  \qquad
  \norm{\tilde s_{k,c}}_0=K_{k,c},
  \label{eq:mask-step-support}
\end{equation}
because \(W_k\) has strictly positive diagonal entries.
The generator is coordinate-sparse when \(\kappa_D=o(D)\), and
constant-coordinate when \(\kappa_D=O(1)\).

\begin{lemma}[Sparse-coordinate scale]
Let \(K_D\) denote a generic draw from the distribution of \(K_{k,c}\).  Then
\begin{equation}
  K_D\stackrel{d}{=}\max\{1,Z_D\},
  \qquad
  Z_D\sim\Bin(D,p_D),
  \label{eq:mask-dist}
\end{equation}
and
\begin{equation}
  \mathbb{E}_D\!\left[K_D\right] = Dp_D+(1-p_D)^D.
  \label{eq:mask-mean}
\end{equation}
For \(D>20\),
\begin{equation}
  \mathbb{E}_D\!\left[K_D\right]
  =
  20+\left(1-\frac{20}{D}\right)^D
  =
  20+e^{-20}+O(D^{-1}).
  \label{eq:mask-large-d}
\end{equation}
Equivalently,
\[
  \kappa_D=20+e^{-20}+O(D^{-1}),
  \qquad
  \frac{\kappa_D}{D}\to0 .
\]
\end{lemma}

\begin{proof}
Before correction,
\[
  Z_D=\sum_{j=1}^{D}\bar b_{k,c,j}\sim\Bin(D,p_D).
\]
The correction changes the cardinality only on \(\{Z_D=0\}\),
where it sets the cardinality to one.  Therefore \(K_D=\max\{1,Z_D\}\), and
\[
  \mathbb{E}_D\!\left[K_D\right]
  =
  \mathbb{E}_D\!\left[Z_D\right]
  +
  \mathbb{P}_D\!\left\{Z_D=0\right\}
  =
  Dp_D+(1-p_D)^D.
\]
Substituting \(p_D=\frac{20}{D}\) gives \eqref{eq:mask-large-d}.
\end{proof}

ENN-surrogate \TuRBO{} uses the equivalent support-size representation: draw
\(K_D=\max\{1,Z_D\}\), then choose \(K_D\) coordinates uniformly without
replacement.  The support-cardinality law is the same as in
\eqref{eq:corrected-mask}.

Thus the candidate generator has active-coordinate scale
\(\kappa_D=20+e^{-20}+O(D^{-1})=o(D)\).  The next section shows that the
failure tolerance has scale at least \(D\).

\section{Failure Counter and Shrink Rule}

Now separate the trust-region update rule from the particular surrogate.  The
surrogate and acquisition function decide which evaluations are successes; the
following counting argument only uses the resulting success/failure sequence.
Fix one region.  Let
\[
  C_k\in\N_0,
  \qquad
  a_k\in\N,
  \qquad
  \eta_k\in\{0,\ldots,a_k\},
  \qquad
  I_k\in\{0,1\},
\]
where \(C_k\) is the failure counter before update \(k\), \(a_k\) is the number
of evaluations assigned to the region at update \(k\), \(\eta_k\) is the
evaluation-normalized failure increment used by the implementation, and
\(I_k=1\) denotes a successful update.  Define the pre-shrink counter
\begin{equation}
  C_{k+1}^{-}
  =
  (1-I_k)(C_k+\eta_k).
  \label{eq:pre-shrink-counter}
\end{equation}
Given a failure tolerance \(\tau_D\), the shrink indicator is
\begin{equation}
  H_k
  =
  \indic{C_{k+1}^{-}\ge\tau_D}.
  \label{eq:shrink-indicator}
\end{equation}
For a shrink factor \(\gamma_{\downarrow}\in(0,1)\), the length update is
\begin{equation}
  \Delta_{k+1}=\gamma_{\downarrow}^{H_k}\Delta_k .
  \label{eq:shrink-update}
\end{equation}
Thus \(H_k=1\) shrinks the region and \(H_k=0\) leaves \(\Delta_k\) unchanged.
The counter reset is
\begin{equation}
  C_{k+1}=(1-H_k)C_{k+1}^{-}.
  \label{eq:post-shrink-counter}
\end{equation}
Thus \(\tau_D\) is the failure tolerance, measured in evaluations assigned to
the region.  This is the scale to compare against \(\kappa_D\).

\begin{theorem}[Shrink-count bound]
Consider any counter-based trust-region method satisfying
\(\tau_D\ge D\).  Let \(T\) be the number of evaluations assigned to a fixed
region after initialization.  Put
\begin{equation}
  A_0=0,
  \qquad
  A_K=\sum_{k=1}^{K} a_k,
  \qquad
  K(T)=\max\{K\in\N_0:A_K\le T\},
  \qquad
  N_{\downarrow}(T)=\sum_{k=1}^{K(T)}H_k .
  \label{eq:shrink-count}
\end{equation}
Then
\begin{equation}
  N_{\downarrow}(T)
  \le
  \floor*{\frac{T}{\tau_D}}
  \le
  \floor*{\frac{T}{D}} .
  \label{eq:main-bound}
\end{equation}
\end{theorem}

\begin{proof}
Let
\[
  F_T=\sum_{k=1}^{K(T)}(1-I_k)\eta_k .
\]
Then
\[
  F_T
  \le
  \sum_{k=1}^{K(T)}a_k
  =
  A_{K(T)}
  \le
  T .
\]
Every shrink event requires failure increments with total mass at least
\(\tau_D\).  By \eqref{eq:post-shrink-counter}, the counter is reset after each
shrink, so the increments certifying distinct shrink events are disjoint.  Hence
\[
  N_{\downarrow}(T)\tau_D\le F_T\le T.
\]
The first inequality in \eqref{eq:main-bound} follows.  The second follows from
\(\tau_D\ge D\).
\end{proof}

\begin{corollary}[No shrink below ambient budget]
If \(T<D\), then \(N_{\downarrow}(T)=0\).  Thus no failure-driven shrink
can occur, even if every update is unsuccessful.
\end{corollary}

\section{\TuRBO{} Failure Tolerance}

We now identify \(\tau_D\) for the reference \TuRBO{} rules.  This is the only
place where implementation details enter the argument.  The official code uses
the minimization convention: a batch is successful if it improves the local
incumbent by the implementation tolerance
\begin{equation}
  y_k^{\min}=\min_{1\le i\le q} f(x_{k,i}),
  \qquad
  I_k=1
  \quad\Longleftrightarrow\quad
  y_k^{\min}
  <
  \best
  -
  10^{-3}\abs{\best}.
  \label{eq:success}
\end{equation}
The exact success criterion only determines when the counter resets; the bound
above holds for every success/failure sequence.

\paragraph{\TurboOne.}
With batch size \(q\), the reference single-region code uses the batch-level
failure tolerance \cite{turboOne}
\begin{equation}
  \widehat{\tau}^{(1)}_D(q)
  =
  \ceil*{
    \max\left\{\frac{4}{q},\frac{D}{q}\right\}
  } .
  \label{eq:turbo-one-failure-tolerance}
\end{equation}
In evaluation units this is
\begin{equation}
  \tau_D^{(1)}(q)
  =
  q\widehat{\tau}^{(1)}_D(q)
  \ge D
  \label{eq:turbo-one-evaluation-threshold}
\end{equation}
Thus one shrink requires at least \(\tau_D^{(1)}(q)\) evaluations assigned
through unsuccessful batches.

\paragraph{ENN-surrogate \TuRBO.}
In ENN-surrogate \TuRBO{} with batch size \(q\), the trust-region failure
clock is batch-level:
\begin{equation}
  \widehat{\tau}^{(\mathrm{ENN})}_D(q)
  =
  \ceil*{
    \max\left\{\frac{4}{q},\frac{D}{q}\right\}
  } .
  \label{eq:turbo-enn-failure-tolerance}
\end{equation}
Under full-batch updates, each unsuccessful update assigns \(q\) evaluations to
the region.  The evaluation threshold is
\begin{equation}
  \tau_D^{(\mathrm{ENN})}(q)
  =
  q\widehat{\tau}^{(\mathrm{ENN})}_D(q)
  \ge D .
  \label{eq:turbo-enn-evaluation-threshold}
\end{equation}
The shrink-count bound therefore applies unchanged to ENN-surrogate \TuRBO{} in
the full-batch regime.

\paragraph{\TurboM.}
The reference multi-region code counts failures by evaluations assigned to each
region and uses \cite{turboM}
\begin{equation}
  \tau_D^{(m)}=\max\{5,D\}.
  \label{eq:turbo-m-failure-tolerance}
\end{equation}
Therefore \(\tau_D^{(m)}=D\) for \(D\ge5\).

\begin{corollary}[\TuRBO{} shrink bound]
For \TurboOne{}, ENN-surrogate \TuRBO{} under full-batch updates, and
\TurboM{}, in dimensions \(D\ge5\), a trust region receiving \(T\)
post-initialization evaluations satisfies
\begin{equation}
  N_{\downarrow}(T)
  \le
  \floor*{\frac{T}{D}} .
  \label{eq:turbo-bound}
\end{equation}
\end{corollary}

\section{Consequences}

The reference rule shrinks by \(\gamma_{\downarrow}=\frac{1}{2}\).  If a region
starts at \(\Delta_0\) and restarts once \(\Delta_k<\Delta_{\min}\), then the
number of shrink events required for one restart is
\begin{equation}
  h_{\mathrm{restart}}
  =
  \min\{h\in\N:\Delta_0 2^{-h}<\Delta_{\min}\}
  =
  \floor*{\log_2\!\left(\frac{\Delta_0}{\Delta_{\min}}\right)}+1.
  \label{eq:restart-height}
\end{equation}
With the reference defaults \(\Delta_0=0.8\) and \(\Delta_{\min}=2^{-7}\),
\[
  h_{\mathrm{restart}}=7,
  \qquad
  T_{\mathrm{restart}}
  \ge
  h_{\mathrm{restart}}\tau_D
  \ge
  7D .
\]
Thus a default restart requires at least \(7D\) evaluations assigned through
unsuccessful updates in the same region.

The mathematical point is the scale separation
\begin{equation}
  \kappa_D
  =
  \mathbb{E}_D\!\left[K_{k,c}\right]
  =
  20+e^{-20}+O(D^{-1}),
  \qquad
  \tau_D\ge D.
  \label{eq:scale-separation}
\end{equation}
Equivalently,
\begin{equation}
  \frac{\kappa_D}{D}\to0,
  \qquad
  \frac{\tau_D}{D}\ge1,
  \qquad
  \frac{\tau_D}{\kappa_D}\to\infty .
  \label{eq:normalized-separation}
\end{equation}
The proposed candidate step is low-support in ambient coordinates; the realized
step after clipping has no larger support by \eqref{eq:clipping-support}.  The
length update, however, waits for an ambient-dimensional number of assigned
evaluations in unsuccessful updates.  Hence the failure-driven length update
cannot be triggered whenever the per-region budget \(T\) is below \(D\).  This
is not a claim that \TuRBO{} is ineffective in all high-dimensional problems;
it is a statement about the ambient-dimensional failure tolerance in the regime
\(T<D\).

\begin{example}
If \(D=100{,}000\) and a region receives \(T=10{,}000\) evaluations, then
\(N_{\downarrow}(T)=0\) under the reference rule, even along an all-failure
trajectory.  A default restart requires on the order of \(700{,}000\) assigned
evaluations in unsuccessful updates in that same region.
\end{example}

\begin{remark}
The ENN surrogate changes the local model and candidate ranking, but not the two
scales isolated above: the RAASP support remains constant-coordinate, and under
full-batch updates the failure-driven length update remains ambient-dimensional.
\end{remark}

\begin{thebibliography}{9}
\raggedright
\small

\bibitem{conn2000trust}
Andrew R. Conn, Nicholas I. M. Gould, and Philippe L. Toint.
\newblock \emph{Trust Region Methods}.
\newblock SIAM, 2000.

\bibitem{eriksson2019turbo}
David Eriksson, Michael Pearce, Jacob R. Gardner, Ryan Turner, and Matthias
Poloczek.
\newblock Scalable global optimization via local Bayesian optimization.
\newblock In \emph{Advances in Neural Information Processing Systems}, 2019.
\newblock \url{https://papers.nips.cc/paper/8788-scalable-global-optimization-via-local-bayesian-optimization}

\bibitem{turboCode}
David Eriksson, Michael Pearce, Jacob R. Gardner, Ryan Turner, and Matthias
Poloczek.
\newblock \TuRBO: Trust Region Bayesian Optimization, official code release.
\newblock \url{https://github.com/uber-research/TuRBO}

\bibitem{turboOne}
David Eriksson, Michael Pearce, Jacob R. Gardner, Ryan Turner, and Matthias
Poloczek.
\newblock Reference implementation of \TurboOne.
\newblock \url{https://github.com/uber-research/TuRBO/blob/master/turbo/turbo_1.py}

\bibitem{turboM}
David Eriksson, Michael Pearce, Jacob R. Gardner, Ryan Turner, and Matthias
Poloczek.
\newblock Reference implementation of \TurboM.
\newblock \url{https://github.com/uber-research/TuRBO/blob/master/turbo/turbo_m.py}

\end{thebibliography}

\end{document}
